\documentclass[12pt,a4paper]{article}
\usepackage[a4paper,margin=18mm]{geometry}
\usepackage[T2A]{fontenc}
\usepackage[utf8]{inputenc}
\usepackage[russian]{babel}
\usepackage{amsmath,amssymb,mathtools}
\usepackage{array,booktabs,longtable}
\usepackage{xcolor}
\usepackage{hyperref}
\usepackage{tikz}

\hypersetup{
  colorlinks=true,
  linkcolor=blue!55!black,
  urlcolor=blue!55!black
}
\setlength{\parindent}{0pt}
\setlength{\parskip}{0.55em}
\renewcommand{\arraystretch}{1.2}

\begin{document}

\begin{titlepage}
\centering
\vspace*{0.22\textheight}
{\LARGE\bfseries Ревью домашней работы\par}
\vspace{2.2cm}
{\large Автор: Лёвин Артём Александрович\par}
\vspace{0.8cm}
{\large 9 сентября 2026 года\par}
\vfill
\begin{tikzpicture}
  \draw[blue!40!black,line width=0.8pt]
    (0,0) .. controls (2.2,0.55) and (4.2,-0.55) .. (6.5,0)
          .. controls (8.2,0.4) and (9.6,0.25) .. (11,0.05);
\end{tikzpicture}
\end{titlepage}

\newpage
\section*{Сводная таблица}
\small
\begin{longtable}{@{}>{\centering\arraybackslash}p{0.07\linewidth}p{0.16\linewidth}p{0.27\linewidth}p{0.16\linewidth}p{0.16\linewidth}@{}}
\toprule
№ задания & Тема & Суть ошибки & Ваш ответ & Правильный ответ \\
\midrule
\endfirsthead
\toprule
№ задания & Тема & Суть ошибки & Ваш ответ & Правильный ответ \\
\midrule
\endhead
\hyperref[task:8]{8} & Квадратное уравнение; дискриминант; извлечение квадратного корня & Верный итог получен при математически некорректной записи квадрата отрицательного коэффициента и некорректной цепочке для оценки $\sqrt{D}$. Требуемый метод нужно записать строго. & $4$ и $76$ & $4$ и $76$ \\
\addlinespace
\hyperref[task:9]{9} & Квадратное уравнение; дискриминант, близкий к $b^2$ & Решение остановлено после записи $D=5002^2-40000$: не найден $\sqrt D$ методом близости к $|b|$, корни не вычислены. Также квадрат отрицательного коэффициента записан без скобок. & не указан & $\dfrac1{50}$ и $50$ \\
\bottomrule
\end{longtable}
\normalsize

\newpage
\thispagestyle{empty}
\vspace*{\fill}
\begin{center}
{\LARGE\bfseries Разбор ошибок}
\end{center}
\vspace*{\fill}

\newpage
\phantomsection\label{task:8}
\section*{Задание 8}

\subsection*{Текст задания}
Решите квадратное уравнение
\[
x^2-80x+304=0.
\]
Корень из дискриминанта найдите методом ближайших квадратов и последней цифры.

\subsection*{Ваше решение}
В работе записано:
\[
D=-80^2-4\cdot304=5184,
\]
после чего для корня из дискриминанта приведена цепочка вида
\[
70^2<\sqrt{5184}\le 80\quad\Longrightarrow\quad 72^2,
\]
а затем получены корни
\[
x_1=\frac{80+72}{2}=76,\qquad
x_2=\frac{80-72}{2}=4.
\]

\subsection*{Ваш ответ}
\[
4\text{ и }76.
\]

\subsection*{Ошибка}
\textcolor{red}{Ошибка:} неверная математическая запись при вычислении дискриминанта и при обосновании значения $\sqrt{D}$.

\subsection*{Где именно возникает ошибка}
Уравнение переписано правильно, коэффициенты можно определить как
\[
a=1,\qquad b=-80,\qquad c=304.
\]
Это последний полностью корректный этап.

Первый неверно записанный шаг:
\[
D=-80^2-4\cdot304.
\]
Если коэффициент $b=-80$, то в формуле дискриминанта требуется квадрат \emph{всего} числа $-80$:
\[
b^2=(-80)^2.
\]
Без скобок выражение $-80^2$ по обычному порядку действий означает
\[
-(80^2)=-6400,
\]
а не $6400$.

Есть и вторая проблема: цепочка
\[
70^2<\sqrt{5184}\le 80
\]
не является корректным сравнением. Слева стоит число $70^2=4900$, а в середине $\sqrt{5184}=72$; поэтому неравенство $4900<72$ ложно. Если сравниваются квадраты, нужно сравнивать $D$ с квадратами. Если сравнивается $\sqrt D$, нужно сравнивать сам корень с обычными числами.

\subsection*{Почему это неверно}
Формула дискриминанта имеет вид
\[
D=b^2-4ac.
\]
Подстановка отрицательного коэффициента требует скобок:
\[
D=(-80)^2-4\cdot1\cdot304.
\]
Скобки здесь не декоративны: они показывают, что возводится в квадрат число $-80$ целиком.

Далее требуется именно метод ближайших квадратов и последней цифры. Корректная логика такая. Сначала локализуем $\sqrt{5184}$ между десятками:
\[
70^2=4900<5184<6400=80^2,
\]
следовательно,
\[
70<\sqrt{5184}<80.
\]
Последняя цифра числа $5184$ равна $4$. Квадрат целого числа оканчивается на $4$, если само число оканчивается на $2$ или на $8$. В интервале от $70$ до $80$ поэтому остаются два кандидата: $72$ и $78$.

Чтобы выбрать между ними, достаточно сравнить с ближайшим удобным квадратом:
\[
75^2=5625.
\]
Так как
\[
5184<5625,
\]
то
\[
\sqrt{5184}<75.
\]
Из кандидатов $72$ и $78$ подходит только $72$. Значит,
\[
\sqrt{5184}=72.
\]

Таким образом, численный ответ в работе получен верно, но требуемое обоснование записано некорректно. Для задания, где способ нахождения $\sqrt D$ указан отдельно, это существенный элемент решения.

\subsection*{Ключевая идея}
\textcolor{blue!60!black}{Ключевая идея:} при отрицательном $b$ обязательно писать $b^2=(-80)^2$, а при поиске $\sqrt D$ сначала ограничить корень соседними квадратами, затем использовать последнюю цифру и дополнительное сравнение, которое однозначно выбирает кандидата.

\subsection*{Исправление от последнего правильного шага}
Берём
\[
a=1,\qquad b=-80,\qquad c=304.
\]
Тогда
\[
D=(-80)^2-4\cdot1\cdot304
=6400-1216
=5184.
\]

Теперь находим $\sqrt D$ требуемым способом:
\[
70^2=4900<5184<6400=80^2,
\]
поэтому
\[
70<\sqrt{5184}<80.
\]
Число $5184$ оканчивается на $4$, значит возможная последняя цифра корня --- $2$ или $8$. Получаем кандидаты $72$ и $78$.

Поскольку
\[
5184<75^2=5625,
\]
имеем $\sqrt{5184}<75$, поэтому
\[
\sqrt{5184}=72.
\]

\subsection*{Полное правильное решение}
Решаем
\[
x^2-80x+304=0.
\]
Коэффициенты:
\[
a=1,\qquad b=-80,\qquad c=304.
\]

Дискриминант:
\[
D=b^2-4ac=(-80)^2-4\cdot1\cdot304=6400-1216=5184.
\]

Найдём $\sqrt{5184}$ без разложения на простые множители. Имеем
\[
70^2=4900<5184<6400=80^2,
\]
то есть
\[
70<\sqrt{5184}<80.
\]
Последняя цифра $D$ равна $4$, поэтому целый квадратный корень может оканчиваться на $2$ или $8$. В указанном промежутке это $72$ или $78$. Далее
\[
5184<75^2=5625,
\]
следовательно, корень меньше $75$. Поэтому
\[
\sqrt D=72.
\]

По формуле корней квадратного уравнения:
\[
x_{1,2}=\frac{-b\pm\sqrt D}{2a}
=\frac{80\pm72}{2}.
\]
Отсюда
\[
x_1=\frac{80+72}{2}=\frac{152}{2}=76,
\]
\[
x_2=\frac{80-72}{2}=\frac{8}{2}=4.
\]

\subsection*{Проверка}
По теореме Виета для уравнения
\[
x^2-80x+304=0
\]
сумма корней должна быть равна $80$, а произведение --- $304$.
\[
76+4=80,
\]
\[
76\cdot4=304.
\]
Оба условия выполняются.

\subsection*{Итог}
\textcolor{teal!60!black}{Итог:} корни $4$ и $76$ найдены верно. Исправить нужно именно запись дискриминанта и строго оформить требуемый способ нахождения $\sqrt{5184}$.

\subsection*{Задача на закрепление}
Решите уравнение
\[
x^2-86x+480=0.
\]
Корень из дискриминанта найдите методом ближайших квадратов и последней цифры.

\newpage
\phantomsection\label{task:9}
\section*{Задание 9}

\subsection*{Текст задания}
Решите уравнение
\[
100x^2-5002x+100=0,
\]
не раскладывая дискриминант на простые множители. Используйте близость $\sqrt D$ к $|b|$.

\subsection*{Ваше решение}
В работе записано:
\[
100x^2-5002x+100=0,
\]
затем начато вычисление дискриминанта:
\[
D=-5002^2-4\cdot100\cdot100
=5002^2-40000,
\]
после чего решение остановлено знаком вопроса.

\subsection*{Ваш ответ}
не указан.

\subsection*{Ошибка}
\textcolor{red}{Ошибка:} решение не завершено; кроме того, квадрат отрицательного коэффициента $b=-5002$ первоначально записан без скобок.

\subsection*{Где именно возникает ошибка}
Правильно определены коэффициенты уравнения:
\[
a=100,\qquad b=-5002,\qquad c=100.
\]
Правильной и полезной является также идея привести дискриминант к виду
\[
D=5002^2-40000.
\]
Это последний содержательно корректный шаг, который можно сохранить.

До него, однако, имеется неточная запись
\[
-5002^2.
\]
Для коэффициента $b=-5002$ должно быть
\[
(-5002)^2.
\]

Главная незавершённость начинается после выражения
\[
D=5002^2-40000:
\]
не найдено $\sqrt D$, не использована указанная в условии близость к $|b|=5002$, и не вычислены корни уравнения.

\subsection*{Почему это неверно}
Здесь не нужно вычислять огромное число $5002^2$, затем вычитать $40000$ и отдельно искать квадратный корень. Условие специально подсказывает более короткий школьный путь.

Поскольку
\[
D=5002^2-40000,
\]
а $40000$ сравнительно мало по отношению к $5002^2$, число $\sqrt D$ должно быть немного меньше $5002$. Ищем целое число вида
\[
5002-k.
\]

Удобно использовать разность квадратов. Если предположить, что
\[
\sqrt D=4998,
\]
то разность между квадратами $5002^2$ и $4998^2$ равна
\[
5002^2-4998^2=(5002-4998)(5002+4998)=4\cdot10000=40000.
\]
Именно столько вычитается в дискриминанте. Следовательно,
\[
5002^2-40000=4998^2,
\]
то есть
\[
\sqrt D=4998.
\]

Это и есть использование близости $\sqrt D$ к $|b|$: вместо прямого вычисления большого дискриминанта мы замечаем, что искомый корень отличается от $5002$ всего на $4$.

\subsection*{Ключевая идея}
\textcolor{blue!60!black}{Ключевая идея:} представить дискриминант как квадрат $|b|$ минус небольшая величина и подобрать близкий квадрат через формулу разности квадратов
\[
A^2-B^2=(A-B)(A+B).
\]

\subsection*{Исправление от последнего правильного шага}
Продолжаем с
\[
D=5002^2-40000.
\]
Замечаем:
\[
40000=4\cdot10000
=(5002-4998)(5002+4998).
\]
Поэтому
\[
5002^2-4998^2=40000,
\]
и значит
\[
D=4998^2.
\]
Следовательно,
\[
\sqrt D=4998.
\]

Теперь применяем формулу корней:
\[
x_{1,2}=\frac{-b\pm\sqrt D}{2a}
=\frac{5002\pm4998}{200}.
\]
Тогда
\[
x_1=\frac{5002+4998}{200}
=\frac{10000}{200}=50,
\]
\[
x_2=\frac{5002-4998}{200}
=\frac{4}{200}=\frac1{50}.
\]

\subsection*{Полное правильное решение}
Для уравнения
\[
100x^2-5002x+100=0
\]
имеем
\[
a=100,\qquad b=-5002,\qquad c=100.
\]
Дискриминант:
\[
D=b^2-4ac
=(-5002)^2-4\cdot100\cdot100
=5002^2-40000.
\]

Не вычисляем $5002^2$ напрямую. По условию используем близость $\sqrt D$ к $|b|=5002$. Проверим число $4998$, которое меньше $5002$ на $4$:
\[
5002^2-4998^2
=(5002-4998)(5002+4998)
=4\cdot10000
=40000.
\]
Следовательно,
\[
5002^2-40000=4998^2,
\]
то есть
\[
D=4998^2,
\qquad
\sqrt D=4998.
\]

По формуле корней:
\[
x_{1,2}=\frac{5002\pm4998}{200}.
\]
Получаем
\[
x_1=50,
\qquad
x_2=\frac1{50}.
\]

\subsection*{Проверка}
Разделим исходное уравнение на $100$:
\[
x^2-\frac{2501}{50}x+1=0.
\]
По теореме Виета произведение корней должно быть равно $1$:
\[
50\cdot\frac1{50}=1.
\]
Сумма корней должна быть равна
\[
\frac{2501}{50}.
\]
Проверяем:
\[
50+\frac1{50}
=\frac{2500}{50}+\frac1{50}
=\frac{2501}{50}.
\]
Оба условия выполняются.

\subsection*{Итог}
\textcolor{teal!60!black}{Итог:} после записи $D=5002^2-40000$ нужно было использовать разность близких квадратов и получить $\sqrt D=4998$. Корни уравнения:
\[
\boxed{50\text{ и }\frac1{50}}.
\]

\subsection*{Задача на закрепление}
Решите уравнение
\[
25x^2-626x+25=0,
\]
не раскладывая дискриминант на простые множители. Используйте близость $\sqrt D$ к $|b|$.

\end{document}
